\paragraph{Sharpness also for the primitive sextic fields.}
The same set of 66 primes is exact for local compatibility in the $A_5$
family. Use the same union $S$ of splitting primes over all $66$ candidates. We construct an $A_5$-extension in
which every $q\in S$ is unramified with Frobenius a $3$-cycle in the
natural quintic action. In the sextic action its cycle lengths are $3,3$.
The proof of local compatibility above then applies unchanged.

Let $X$ parametrize monic quintics $f$ together with $y\ne0$ satisfying
$y^2=\operatorname{disc}(f)$. Its function field is
$\mathbb Q(x_1,\ldots,x_5)^{A_5}$: the elementary symmetric functions of
the roots generate the $S_5$-invariant field and the Vandermonde product
generates its quadratic extension. Maeda proved this $A_5$-invariant
field is purely transcendental. Consequently $X$ has weak approximation
at finite sets of primes. Indeed, a nonempty Zariski open of $X$ is
isomorphic to an open in affine $5$-space. Every nonempty local open on
$X(\mathbb Q_q)$ meets this chart since $X$ is smooth and geometrically
integral. Ordinary weak approximation in the five coordinates now
suffices. See \href{https://doi.org/10.1016/0021-8693(89)90174-9}
{T. Maeda, \emph{Noether's problem for $A_5$}, J. Algebra 125 (1989)}.

At every $q\in S$ prescribe squarefree reduction of type $(3,1,1)$.
Such a polynomial exists even over $\mathbb F_2$, by taking an irreducible
cubic and two distinct linear factors. At two further primes prescribe
types $(3,1,1)$ and $(5)$. Each prescription defines a nonempty local
open on $X$. In fact, a monic integral lift of the chosen reduction has
unramified splitting field with Frobenius given by the indicated cycle
type. This permutation is even, so the signed Vandermonde product of its
roots belongs to $\mathbb Q_q$ and supplies $y$. Weak approximation now
gives one rational quintic realizing all prescriptions. Its group is
contained in $A_5$ and has order divisible by $15$. A subgroup of order
$30$ would contradict simplicity. A subgroup of order $15$ would have
normal Sylow $5$-subgroup, whereas the normalizer of this subgroup in
$A_5$ has order $10$. Hence the splitting field has group $A_5$, as required.
